% Source: sources/returns/S6_04_MAPS_MANIFOLDS/sections/manifolds.tex:1150–1312
\paragraph{Theorem D.2.6 (Proved: infinitely many orders for $III^*+3I_1$).}
Use the corrected $III^*,I_1,I_1,I_1$ construction of K11, with
two primitive rank-one Mumford fillings and one rank-two hexagonal
filling. In the ordinary homology basis $(e_0,e_1,e_2,e_3)$ dual
to $(\psi,u,w,\delta)$, put $\ell=e_0-e_1-e_2$.
At $III^*$ one may make a logarithmic transformation with
\[
 m=4k,\qquad v=n\ell+r e_3,\qquad
 k>0,\quad n,r\in\mathbb Z,\quad \gcd(n,4k)=1.
\]
For every such datum the resulting compact complex threefold has
\[
 \pi_1(X)=\mathbb Z/|n|,\qquad
 H_0=H_6=\mathbb Z,\quad H_5=0,\quad
 H_1=H_2=H_3=H_4=\mathbb Z/|n|.
\]
Consequently $n=1$ or $n=-1$ gives a complex structure on $S^6$
for every $k>0$, with no coprimality restriction on $k$.
Thus this configuration answers Viaclovsky's Question 4
affirmatively. This assertion concerns the displayed analytic
fillings and integer twists, not all modifications of that surface.
\begin{proof}
The geometric input is K11, checked against
\cite[Proposition 3.1, Lemma 3.2, Proposition 4.2]{manViaclovskyTwo26}.
The order-four good-reduction family is smooth and proper and its
punctured comparison preserves the distinguished section.
Pull it back by a degree-$k$ map of the covering disc, so that the
original base coordinate is $s^{4k}$. The linear deck action still
has order four on the torus lattice. It fixes $\ell$ and $e_3$.
Translation by $v/(4k)$, combined with rotation of $s$, defines
an order-$4k$ action. On the central torus, the invariant map
$\psi:T^4\to\mathbb R/\mathbb Z$ changes by $n/(4k)$;
coprimality proves freeness. The punctured comparison is changed
by the holomorphic translation
$[\log(s)v(s)/(2\pi i)]$, exactly as in D.2.5.
The quotient is therefore smooth and glues to the unchanged
three Mumford pieces. Properness follows from properness of the
four fillings and compactness of the base. Its meridian relation
is $g^m=v$. This constructs the claimed analytic threefold for
every displayed parameter, rather than extrapolating a finite table.

Here is the integral topology. The cohomological matrices, in the
stated marking, are
\[
 \begin{aligned}
 T_0&=\begin{pmatrix}1&0&-2&-1\\0&0&-1&0\\0&1&0&-1\\0&0&0&1\end{pmatrix},
 &T_1&=\begin{pmatrix}1&0&2&2\\0&1&1&1\\0&0&1&0\\0&0&0&1\end{pmatrix},\\
 T_2&=\begin{pmatrix}1&0&0&0\\0&1&0&0\\0&-1&1&0\\0&0&0&1\end{pmatrix},
 &T_3&=\begin{pmatrix}1&0&0&-1\\0&1&1&0\\0&0&1&0\\0&0&0&1\end{pmatrix}.
 \end{aligned}
\]
Let $A_i=T_i^{-t}$. The primitive vanishing subgroups of the
three Mumford pieces are respectively
$\mathbb Z(e_2+e_3)$, $\mathbb Ze_1$, and
$\mathbb Ze_2\oplus\mathbb Ze_3$.
They kill $e_1,e_2,e_3$. Their meridians are killed by the
extending section, and the sphere relation then kills the finite
meridian as well. The finite relation $g^m=v$ leaves exactly
$ne_0=0$. All boundary fundamental-group maps are onto, so
no generators have been omitted. This proves $\pi_1(X)=\mathbb Z/|n|$.

We justify the rank-one normal tori used in the degree-two
calculation. In a primitive rank-one Mumford filling, the
normalization is a flat $\mathbb P^1$-bundle over an elliptic
curve, with its two sections identified by a translation.
Topologically trivialize the flat line bundle and isotope the
translation to the identity, keeping the extending section fixed.
The collar pair is then the product of the primitive nodal
elliptic filling with $T^2$; its local equation is $xy=q$.
It follows that $H^3$ of its central fibre is free of rank three.
The normal tori are the sweeps of all three integral invariant
loops: the vanishing circle, whose rotation extends over $xy=q$,
and the two loops of the remaining elliptic factor. They are
boundaries of relative three-chains obtained from the section
disc and those rotations or translations. The unit-factor
deformation leaves the marked boundary and section comparison
unchanged up to isotopy. For the rank-two filling, the same
argument gives the two invariant angular loops and $H^3=\mathbb Z^2$.
Thus these are integral normal-torus classes; retaining only the
vanishing loop in a rank-one piece would omit two relations.

For completeness we give primitive-minor certificates for the
global lattice calculation. Set
\[
 D=\bigl[(A_0-I)\ (A_1-I)\ (A_2-I)\bigr],
 \qquad
 j_3(x)=-(A_1A_2x,A_2x,x).
\]
For $i=0,1,2$ let $j_i(x)$ place $x$ in block $i$ of
$\Lambda^3$. The fixed lattices at the three Mumford points have
the respective integral bases
\[
 (-e_0+2e_1,e_2,e_3),\qquad (e_0,e_1,e_3),\qquad(e_2,e_3).
\]
Form a $12\times9$ matrix $K$ whose first eight columns are
their images under $j_1,j_2,j_3$, in that order, and whose ninth
column is $j_0(\ell)$. Direct multiplication gives $DK=0$ and
$\operatorname{rank}D=3$. The minor of $K$ using rows
\[
 (1,2,4,5,7,8,9,10,12)
\]
is $-1$, with rows numbered from one. Hence its columns form
an integral basis of $\ker D$. Moreover
$j_0(e_3)=-K_3-K_6-K_8$, so the finite $e_3$ sweep vanishes
after the Mumford relations.

In the exterior basis
$(e_0e_1,e_0e_2,e_0e_3,e_1e_2,e_1e_3,e_2e_3)$,
the matrix $[\bigwedge^2A_i-I]_{i=0}^3$ has primitive rank five:
rows $(1,2,3,5,6)$ and columns $(1,2,4,5,7)$ have minor $-1$.
Its cokernel is detected by
$\xi=u\wedge w+2\psi\wedge\delta$, whose value on $e_1e_2$
is one. Thus the torus-bundle Serre sequence gives
$H_2(V)=\mathbb Z^{10}$, generated by this fibre class and
the nine sweep classes defining $K$. Each sweep has zero
integral against the flat form $\xi$; these classes consequently
fix an integral splitting, not only a rational one.
The eight Mumford normal tori span a primitive subgroup.
Their domain has rank eight, and the preceding local construction
realizes all eight columns, so their image is exactly that subgroup.
The quotient is
\[
 Q=\mathbb Z\langle s,t\rangle,
 \qquad s=[e_1\wedge e_2],\quad t=[T_{\ell,0}].
\]

Abelianization of the finite fibre group gives
$e_2=-e_1$, $e_3=2e_1$ and
$mg=ne_0+2r e_1$. Since $\gcd(n,m)=1$, it is free of rank
two with basis $e_1,z$, where
$z=\alpha e_0+\beta g$ and $m\alpha+n\beta=1$.
The normal-circle formulas proved in D.2.5 apply. Here
\[
 \mathcal N_me_1=2k e_3,\qquad
 \mathcal N_me_0=m\ell,\qquad
 \xi(e_1,v)=-n,\qquad \xi(e_0,v)=2r.
\]
The first norm dies in $Q$. The normal tori of $e_1,e_0,g$
therefore have coordinates $(n,0)$, $(-2r,m)$ and $(0,n)$.
In the basis $e_1,z$, the remaining attachment matrix is
\[
 \begin{pmatrix}n&-2r\alpha\\0&1\end{pmatrix}.
\]
It has Smith form $(1,|n|)$ for every $k$ and $r$.

To check that this cokernel is actually $H_2(X)$, the boundary
$H_1$ groups have ranks $3,4,4,3$ and are free; the target
of their Mayer--Vietoris map has rank
$4+2+3+3+2=14$ and is also free. Its cokernel is the finite
group $H_1(X)=\mathbb Z/|n|$, so its kernel is zero.
Every central fibre is irreducible and a multiple of it is
principal; the pair-intersection argument in D.2.5 makes each
$H_2(M_i)\to H_2(N_i)$ onto. Thus $H_2(X)$ is precisely
the displayed normal-torus cokernel. Euler localization gives
$\chi(X)=2$, and integral duality gives all remaining groups
as in D.2.5. Sphere recognition for $|n|=1$ is the same argument.

The script \texttt{code/manifolds\_iii\_star.py} executes all
stated lattice certificates, the norm identities, and 312 actual
$10\times10$ Smith forms for $1\leq k\leq12$,
$-15\leq n\leq15$, $\gcd(n,4k)=1$, and $r=0,3$.
The proof is independent of that finite validation bound.
\end{proof}
